\section{When every decision is preserved}
\label{sec:flat}

A finite condition on the driving prototypes characterizes preservation
of all strict decisions at all finite times. The classifier in this section
is an arbitrary finite affine-score readout with distinct affine functions
and full-dimensional strict decision cells, ignoring functions with empty
strict cells. The driving field is the Gaussian mixture field above.
With one distinct prototype per class this includes the power diagram
of the class scores \cite{aurenhammer}, with equal masses giving ordinary
Voronoi cells. The distinction matters because the later benchmark uses
multiple prototypes per class and in general has curved faces.

For a cell facet, let $H$ be its supporting hyperplane, let $n$ be a
unit normal, and write the driving prototypes as $a_i=o+u_i+r_i n$
with $o\in H$ and $u_i\perp n$. For each distinct tangential projection
$u$, define
\begin{equation}
 A_u=\sum_{i:u_i=u}w_i e^{-r_i^2/(2\beta)},\qquad
 B_u=\sum_{i:u_i=u}w_i r_i e^{-r_i^2/(2\beta)}.
 \label{eq:fiber}
\end{equation}

\begin{lemma}[Tangency from a sampled gradient flow]
\label{lem:sampled-tangency}
Let $F=\nabla U$ be real analytic on $\R^d$, with a complete forward
flow and compact superlevel sets of $U$. If $\Phi_T(H)\subseteq H$ for
an affine hyperplane $H$ and some $T>0$, then $F\cdot n=0$ on $H$.
\end{lemma}

\begin{proof}
Start an orbit at any point of $H$. All samples at times $kT$ stay
in $H$ and in a compact superlevel set of $U$. The potential increases along the
orbit and is bounded above. Thus the potential increments between
successive samples tend to zero. At a subsequential limit $p\in H$,
continuity gives $U(\Phi_T p)=U(p)$. Since the difference equals
$\int_0^T\lVert F(\Phi_s p)\rVert^2\,\dd s$, we have $F(p)=0$.
If $d=1$, then $H=\{p\}$ and this proves the claim.

Write $x=p+y+rn$ with $y$ tangent to $H$. At $p$ the tangent space to
$H$ is invariant under the derivative $D\Phi_T(p)=e^{TDF(p)}$.
Since $DF(p)$ is real symmetric, its eigenvalues are real. A polynomial
interpolating $T^{-1}\log$ on the finite spectrum of $e^{TDF(p)}$
recovers $DF(p)$, which therefore also leaves the tangent space invariant.
By symmetry it leaves the normal space invariant.
Let $A$ be its tangential block and $b$ its normal entry. If $F\cdot n$
on $H$ is not identically zero, its Taylor expansion at $p$ has a
first nonzero homogeneous part $Q_k(y)$ of degree $k\ge2$. For initial
conditions $(y,0)$ the tangential component of the flow has linear
part $e^{tA}y$. In the normal field, all terms except $br$ and the
restriction to $r=0$ contain a factor $r$ times a term of positive
degree. Hence, if the normal flow terms of degrees below $j$ vanish,
these terms cannot contribute at degree $j$. For $j<k$ the remaining
degree-$j$ term of the normal flow solves $\dot r_j=br_j$ with zero
initial data. Induction gives zero at every such degree. At degree $k$
the only additional contribution is $Q_k(e^{tA}y)$, so the term at $T$ is
\[
 \int_0^T e^{b(T-s)}Q_k(e^{sA}y)\,\dd s.
\]
In a real eigenbasis of $A$ with eigenvalues $\lambda_j$, each monomial
$y^m$ in $Q_k$ is multiplied by
$\int_0^T e^{b(T-s)+s\sum_jm_j\lambda_j}\,\dd s>0$.
The integral cannot be the zero polynomial. This contradicts
$\Phi_T(H)\subseteq H$, proving tangency by analyticity of $F$ on $H$.
The argument allows zero eigenvalues. The integrated homogeneous
operator is the one used in the embedding of maps in autonomous
flows \cite[Eq.~(2.13)]{zhang}. The real spectrum rules out the weak
resonances in the corresponding formal embedding result \cite{ribon}.
\end{proof}

\begin{theorem}[Criterion for affine decision cells]
\label{thm:flat}
The following conditions are equivalent.
\begin{enumerate}
\item At some fixed $T>0$, the map $\Phi_T$ sends every strict decision
cell into itself.
\item Every strict decision cell is invariant for all $t\ge0$.
\item For every projection group on every facet hyperplane, $B_u=0$.
\end{enumerate}
\end{theorem}

\begin{proof}
Under the first condition, continuity maps the relative interior of each
facet into the closures of both incident cells, hence into $H$.
The finite-time flow of an analytic field depends analytically on its
initial point \cite{carrillo}, so the normal coordinate of $\Phi_T$
vanishes on the entire plane in coordinates based at $o\in H$.
The Gaussian log-density potential has compact superlevel sets because
it tends to $-\infty$ as $\|x\|\to\infty$. Lemma~\ref{lem:sampled-tangency}
now gives $F\cdot n=0$ on $H$.
For $z=o+y\in H$, the Gaussian normal velocity is
\begin{equation}
 n^\top F(z)=
 \frac{\sum_u B_u e^{-\lVert y-u\rVert^2/(2\beta)}}
      {\sum_u A_u e^{-\lVert y-u\rVert^2/(2\beta)}} .
 \label{eq:facetvelocity}
\end{equation}
After division by the positive common factor $e^{-\lVert y\rVert^2/(2\beta)}$,
the numerator is a finite sum of exponentials with distinct directions.
Restrict to a line segment in the relative interior whose direction
has distinct inner products with the $u$'s. Consecutive derivatives
at an interior point give an invertible Vandermonde system, so each
$B_u$ vanishes. In dimension one there is a single group and the same
conclusion is immediate. Conversely, vanishing coefficients make the
entire supporting plane invariant by ODE uniqueness. A trajectory on
either strict side cannot reach the plane at finite time, by backward
uniqueness from a hypothetical contact. Each open cell is the
intersection of its open facet halfspaces. This gives the second
condition, which implies the first.
\end{proof}

Classical tangency criteria characterize invariant sets in terms of the
vector field \cite{nagumo,brezis}. Theorem~\ref{thm:flat} starts with decision
preservation at a single sampling time, then replaces the resulting
continuum of tangency conditions by finite equalities in the Gaussian
driving measure. The forward and backward invariance also preserves
closed cells and the set of maximizers on their boundaries when all
affine functions have full-dimensional cells.

\paragraph{A single relaxed mean-shift step.}
For $0<\alpha<1$, preservation of every strict cell by
$x\mapsto (1-\alpha)x+\alpha\mu(x)$ is equivalent to the same balances.
Preservation implies that the step maps each relative open facet into
its support plane, which makes the normal numerator in
\eqref{eq:facetvelocity} vanish there. Conversely, for a balanced plane
$n^\top(\mu(z)-o)=0$ on $H$ and $\partial_n(n^\top\mu)\ge0$ everywhere.
The normal coordinate of the step is a convex combination of the
input's signed distance and a quantity of the same weak sign.
This preserves each open halfspace. For $\alpha=1$ the result holds
if the driving prototypes have nonzero spread in the normal direction
of every facet, as they do with one site per class in the power diagram.
Then the derivative is strictly positive by the posterior variance
identity. Thus one ordinary Gaussian mean-shift step, one exponential
Euler step of positive finite length, and the exact flow have the same
condition for global preservation, although they can transport
individual queries differently.

\begin{corollary}[Reflection rigidity]
\label{cor:reflection}
For distinct equal-mass sites with one site per class, the following are
equivalent.
\begin{enumerate}
\item All strict Voronoi decisions are invariant for every bandwidth.
\item Reflection in every Voronoi facet hyperplane permutes the sites.
\item In their affine span, the sites are a single orbit of a finite
orthogonal reflection group after translating its fixed point to the
origin.
\end{enumerate}
These also follow from preservation at one positive horizon at a
single bandwidth $\beta\ge D^2/4$, where $D$ is the site diameter.
\end{corollary}

\begin{proof}
The equal-mass cells are independent of $\beta$. If they are invariant
on any open interval of $\beta$, each $B_u$ in \eqref{eq:fiber} vanishes
on that interval. Grouping equal values of $r_i^2$ and using independence
of exponentials in $1/(2\beta)$ forces equal masses at $r$ and $-r$
for every $r>0$ in each group. This is precisely the stated reflection
symmetry, which also makes every $B_u$ zero.

At $\beta\ge D^2/4$, the covariance identity gives $DF\preceq 0$.
The Gaussian log-density potential is coercive from above and strictly
concave. Indeed, an affine segment on any line would extend by real
analyticity to an affine restriction on the whole line, contradicting
quadratic decay. There is a unique maximum $p$. A forward trajectory
stays in a compact superlevel set, and the potential derivative is
$\lVert F\rVert^2$. Bounded derivatives on that compact set give uniform
continuity of $\lVert F\rVert^2$ in time. Its finite integral implies
$F(\Phi_t x)\to0$, and hence $\Phi_t x\to p$. Invariance would put $p$
in the closure of every cell, so all sites lie on a common sphere
centered at $p$ and every facet plane contains $p$. In each tangential
projection group, a line normal to the plane intersects that sphere
at most at two points with opposite offsets. Since $B_u=0$, every
nonzero offset has its reflected partner.

If the facets are mirrors of the site set, their reflections restrict
to the affine span and fix the centroid. They generate a finite group
because an isometry of the span is determined by its action on the
sites. An adjacency path through Voronoi cells crosses facets, each of
whose reflections swaps the two incident sites. The adjacency graph is
connected, as a generic path avoids intersections of codimension at
least two. The group therefore acts transitively. Conversely, for an
orbit of a finite real reflection group, the sites have a common norm
around the fixed point. Their Voronoi cells are the normal cones of
the orbit polytope. Each such cone is a union of closed reflection
chambers \cite{ardila}. To see this, in the interior of a chamber a
maximizer of the inner product over the orbit lies in that closed
chamber: reflecting any orbit site across a chamber wall it violates
strictly increases the inner product. There is a unique orbit
representative in a closed chamber, including on its walls.
Thus every Voronoi facet is supported on a reflecting hyperplane.
\end{proof}

Reflection symmetry can also preserve some multicenter classifiers.
Take directions $a_c$ on a common sphere forming a reflection-group orbit.
For each class use prototypes $\lambda_j a_c$ with $\lambda_j>0$ and
weights $w_j>0$ common across classes. The class scores depend on $c$
through the same strictly increasing function of $a_c^\top x$ after
removing the common Gaussian factor. Thus each decision cell is still
a Voronoi cell of the $a_c$'s. Its facet reflections preserve the
driving measure at every radius, making the balances zero at every
bandwidth.

\paragraph{Generic crossings and an exceptional cancellation.}
In dimension $d\ge2$, for at least three sites with equal masses, almost
every configuration has a permeable patch on each Voronoi facet at every
$\beta>0$. Outside finitely many proper algebraic sets, no third site
lies on a pair bisector and no other pair difference is parallel to
that pair's difference. Each third site then has a singleton projection
group with a nonzero offset. By Theorem~\ref{thm:flat}, $F\cdot n$ is
nonzero on an open patch of that facet. Continuous dependence gives
an open set of nearby initial decisions that switch across the facet
at sufficiently small horizons. At any fixed positive horizon,
the time map fails to take some point in the relative interior of
each such facet into its supporting plane. Otherwise the proof of
Theorem~\ref{thm:flat} would force the plane to be tangent. Nearby
initial points on at least one side reach the opposite side of the
supporting plane. The image of that open set is open, so some
points have a different strict label. Invariance without
reflection is possible at a fixed bandwidth and unequal weights.
For sites $(-1,0,1)$ take weights $(1,\sqrt{2/5},1)$ and
$\beta=(2\log(5/2))^{-1}$. The active ties are at $\pm1/4$. At $1/4$,
in units of the common leading score, the normal numerator is
$1/2-(5/4)(2/5)=0$. Symmetry about zero makes the other tie stationary.
Neither tie is a mirror of the sites.

\paragraph{A finite-time departure bound.}
For queries on the positive side of $H$, define
$\epsilon_H^+=\max(0,-\min_u B_u/A_u)$. On the negative side use
$\epsilon_H^-=\max(0,\max_u B_u/A_u)$. Equation~\eqref{eq:facetvelocity}
makes the plane's normal velocity a convex average of the group means.
At distance $s>0$ from the plane on the query's side, the inward
velocity toward that side is at least $-\epsilon_H^\pm-s$ because
$\partial_n(n^\top F)=\operatorname{Var}(n^\top a_i)/\beta-1\ge-1$.
Up to first contact, comparison gives
$s(t)\ge s(0)e^{-t}-\epsilon_H^\pm(1-e^{-t})$. Thus leaving a strict
cell by $T$ requires initial distance at most
$\epsilon_H^\pm(e^T-1)$ from at least one bounding plane.
The resulting union of slabs gives a probability bound under any
feature distribution and a query-level preservation test with no ODE
integration. It is an exact-arithmetic bound. Using additional
pairwise planes in the test keeps it valid and can reduce coverage.
The switch probability also bounds the absolute change in expected
zero-one loss under any joint feature and label law.

\paragraph{A bound on the reachable projection.}
The extrema over the entire plane can be conservative when the query's
trajectory only visits a bounded part of it. Orient $n$ toward the
initial query and let $(y_0,s_0)$ be its tangential and normal
coordinates relative to $o$, with $s_0>0$. The tangential trajectory
obeys an equation of the form $\dot y=\sum_i\pi_i u_i-y$. Thus throughout
$[0,T]$ its projection lies in
\begin{equation}
 K_T=y_0+(1-e^{-T})\operatorname{conv}
       \big(\{0\}\cup\{u_i-y_0\}_i\big).
 \label{eq:projected-tube}
\end{equation}
At a projected point $y\in K_T$, let $L_i(y)$ be the log of each
weighted plane contribution relative to a reference site's.
The shared quadratic term cancels, making $L_i$ affine in
$y$. Its lower and upper extrema $\underline L_i,\overline L_i$
come from evaluating it at the generators of the projected tube.
In an equal-mass Voronoi classifier the two tied sites have relative
likelihood one if either is the reference. Their opposite normal
offsets cancel, so their combined relative mass is two with offset
zero.

Set $d_0=s_0/(e^T-1)$ and $k_i=r_i+d_0$, using the combined offset
where appropriate. A sufficient test for preservation on this side is
\begin{equation}
 \sum_{i:k_i\ge0} k_i e^{\underline L_i}
 +\sum_{i:k_i<0} k_i e^{\overline L_i}>0 .
 \label{eq:local-plane}
\end{equation}
The combined tied mass multiplies its term by two.
The left side lower-bounds the numerator of the normal velocity shifted
by $d_0$, throughout the reachable projection of the trajectory.
Monotonicity of the normal posterior mean away from the plane gives
$\dot s>-s-d_0$ up to first contact, which keeps $s(t)>0$ through $T$.
Applying this test to each bounding plane certifies the initial cell.
The global and local tests can be combined. The global test bounds a
convex average over all projected positions, while the local test
uses the reachable tube at the cost of separating the logit extrema.
Both are sufficient conditions in exact arithmetic.

\subsection{Shared anisotropy and the choice of gradient}
\label{sec:anisotropy}
Consider a shared precision matrix $P/\beta$, with $P$ symmetric
positive definite, in place of the spherical kernel. The Euclidean
gradient of the same scaled log marginal density is now
$F_P=P(\mu-x)$, where responsibilities use this precision. It is
complete and analytic, with the symmetric bounded derivative
$P\Cov(a)P/\beta-P$ and compact superlevel sets for its potential.
Consequently the sampled-time argument in Theorem~\ref{thm:flat}
still forces tangency on every facet plane.

\begin{corollary}[Aligned normals under Euclidean anisotropic ascent]
\label{cor:anisotropy}
For $F_P$, global preservation of the strict affine decisions at
one positive time is equivalent to both of the following conditions on
each facet plane: $Pn=\lambda_n n$ and
$\sum_{i:u_i=u}w_i r_i e^{-\lambda_n r_i^2/(2\beta)}=0$ for each
tangential projection group.
\end{corollary}

\begin{proof}
On an invariant plane, $n^\top P x=n^\top P\mu(x)$, which is bounded
because $\mu(x)$ lies in the prototype hull. The linear function
$n^\top P x$ is bounded over the whole plane only if $Pn$ is parallel
to $n$. When $Pn=\lambda_n n$, the Gaussian exponent separates
into a tangential quadratic form and the normal term
$\lambda_n r_i^2$. The same exponential-independence argument
applies on the tangential space because the restriction of $P$ is
invertible there. This yields the stated balances. Conversely these
balances make the field tangent, and uniqueness preserves the cells.
\end{proof}

For one site per class the face normals are proportional to
$P(a_b-a_a)$. If $P$ is not a multiple of the identity, almost every
site pair has a difference that is not an eigenvector, so even a
generic binary readout admits switches at every positive time.
For instance take $P=\operatorname{diag}(2,1)$ and equal sites
$\pm(1,1)$. On their tie plane $2x_1+x_2=0$ the posterior mean is zero.
At $(1,-2)$, the Euclidean field is $(-2,2)$, which has nonzero
normal component. The preconditioned field $\mu-x$, obtained by
applying $P^{-1}$ to the Euclidean gradient, preserves the labels.
Under whitening $x\mapsto P^{1/2}x$ the preconditioned field is
the spherical model with transformed sites. This distinguishes
the choice of ascent metric from the use of a full bandwidth matrix
in mean shift \cite{chaconduong}.
